\section{Related Work}
\label{sec:related}

Five lines of work each address part of what \voss{} needs: untrusted hosts, churn, latency targets, planned trust and warm state, gray-failure detection and verifiable accounting. None covers them together.

\textbf{Volunteer computing.} BOINC-style systems harvest idle machines for scientific batch work~\cite{anderson2002seti,anderson2004boinc}. Because hosts are untrusted, they replicate tasks and accept a result when a quorum agrees, and they use long deadlines to tolerate machines that disappear.
Replication multiplies cost, presumes deterministic tasks without side effects, protects integrity but not the confidentiality of inputs, and works on a time scale of hours. \voss{} instead removes the integrity and confidentiality concern with attestation, so redundancy is needed only against omission, and then only for risky invocations.

\textbf{Serverless scheduling and keep-alive.} The cold-start literature decides when to keep sandboxes alive or pre-warm them~\cite{shahrad2020wild}, treats keep-alive as caching~\cite{fuerst2021faascache}, or speeds up start-up with lighter sandboxes and snapshots~\cite{agache2020firecracker,oakes2018sock,ao2022faasnap,cadden2020seuss}.
Low-latency schedulers place short tasks by sampling~\cite{ousterhout2013sparrow}, and cluster managers pack tasks onto fleets~\cite{verma2015borg}. All assume trusted, effectively always-on machines. A warm sandbox is lost only by a policy decision, never because an owner closed the lid, and there is no resource corresponding to a provisioned secret.
\voss{} keeps the start-up vocabulary but adds a per-session trust step and plans warm state under churn and exposure constraints.

\textbf{Harvested and edge serverless.} Serverless platforms on harvested datacenter VMs handle evictions~\cite{zhang2021harvest}, and edge schedulers place containers near data and users~\cite{rausch2021edge}. Their hosts are run by the provider or a trusted party, so the failure model is eviction, not silent omission by an adversary, and no secret is entrusted to a party that might read it.

\textbf{Confidential serverless.} Trusted execution environments protect code from the host operating system and the platform operator~\cite{arnautov2016scone,costan2016sgx,amd2020snp}, and recent work targets serverless functions in confidential VMs and containers~\cite{wallet2025,segarra2024sesame,misono2025cvm}.
These designs address confidentiality on provider-managed servers. They do not address machines whose sessions end unpredictably, and they do not treat establishing trust as a scheduled resource with a cost, a lifetime and an exposure. \voss{} is substrate-agnostic: the same scheduler runs over a TPM-attested appliance or over confidential VMs.

\textbf{Trusted components and failure detection.} Small trusted components, such as a counter or log, restrict what a faulty node can do and reduce the replication needed for Byzantine agreement~\cite{levin2009trinc,veronese2013minbft}. We use the idea differently: attestation does not make agreement cheaper, it makes the faulty node \emph{non-Byzantine}, which lets us replace agreement among volunteers by a provider-run referee.
Heartbeat and accrual detectors~\cite{hayashibara2004phi} infer liveness, which is not progress. Gray failures~\cite{huang2017gray} motivate watching end-to-end progress, which \voss{} does at the trusted router and the attested scheduler.

\textbf{Hedging and robust provisioning.} Hedged requests reduce tail latency by duplicating slow work~\cite{dean2013tail}. \voss{} hedges in proportion to estimated risk, on a different failure domain, within a budget. For the planner we use budgeted robust optimisation~\cite{bertsimas2004price}, which fits the question ``is demand still covered if any $f$ hosts fail together?'' without a probability model for adversaries.
